\section{The \instname\ data reduction flow.}
\label{sec:reduction_flow}
\input{reduction_flow}


\subsection{Reference spectrum}

This tasks selects the reference spectrum to fit in order to derive
the parameters of the atmosphere (such as molecules) and the
instrument (such as spectral resolution).

The user has the possibility to select whether this has to be done
directly on the input science spectrum to correct, or on a telluric
standard star that is somehow associated to the science spectrum to
correct. This is determined by the workflow parameter {\bf
use\_standard} (default ``true'').


{\bf The current version of the \wkfname\ workflow only support the
association of telluric standard stars for XSHOOTER observations.}

A telluric standard star is recognized as such if it has the string
{\tt TELL\_SLIT} in its product category. It must be in a format
compatible with the 1D spectra of the ESO archive standard.

In order to be associated to the science (only for XSHOOTER
observations) it must match the {\tt SEQ.ARM}, {\tt INS.OPTI4.NAME}
(for VIS observations), and {\tt INS.OPTI.5.NAME} (for NIR
observations). Note: telluric correction of XSHOOTER UVB data is not
supported.

Once the reference spectrum is identified (either telluric standard or
the very same science spectrum to correct), it is passed to the next
task.

\subsection{Parameter configuration}

This task analyzes the input spectrum and configure the parameters
that will be used by the recipe {\tt molecfit\_model} for the fit,
depending on the instrument, observational setup, and wavelength
coverage. Note that this step does nor run any fit, it creates a file
containig the values of the recicpe parameters to be passed to the
next task and used by {\tt molecfit\_model}.

Typically, there is a task per instrument (xshooter has 2, one for the
VIS and one for the NIR arm). But the automatic parameter adjustion is
done separately for each dataset.

By default, the configured parameters are directly passed to the next
steps. If one desires to review the parameters and change before
starting the fit, the workflow parameter {\bf set\_parameters\_only}
has to be set to ``true'' {\it before starting the reduction}. This will be
explained in more details in Section 
\ref{sec:adjust_recipe_parameters}.

\subsection{Atmospheric model}

This task runs the recipe {\bf molecfit\_model} on the reference
spectrum. It uses the recipe parameters defined on the previous task.
Fitting parameters cannot be changed from this task: fine tuning of
the fit must be done by setting again the recipe parameters on the
previous task using the configuration editor (see Section
\ref{sec:configuration_editor}).

\subsection{Atmospheric transmission}

This task runs the recipe {\bf molecfit\_calctrans}. It uses the
atmospheric parameters computed on the reference spectrum and computes
the full atmospheric transmission on the wavelenght range defined by
the scientific spectrum.

\subsection{Telluric correction}

This task applies the atmospheric transmission correction to the
scientific spectrum.
